% GENERATED FILE. Edit canonical lessons, then run npm run generate:books.
% canonical-source-hash: fnv1a64:36bde077aa3a3206
% canonical-lessons: MR-C427-vyaj, MR-C427-thakbaki, MR-C427-udhar, MR-C427-baki, MR-C427-phedne, MR-R427-review-420-427

\chapter{Interest and Arrears}
\label{ch:mr-interest-and-arrears}
\hlchaptermodality{427}

\begin{chapteropening}
\textbf{By the end of this chapter:} \emph{I can talk about interest, arrears, buying on credit, what is left, and repaying.}

Complete the last lesson of chapter 427: I can talk about interest, arrears, buying on credit, what is left, and repaying.
\end{chapteropening}

\section[\texorpdfstring{\mr{व्याज} (vyāj)}{व्याज (vyāj)}]{\mr{व्याज} (vyāj) — interest on money}

\label{lesson:MR-C427-vyaj}



\begin{quote}
Before the new one: say the Marathi for a sum of money, then the Marathi for to deposit.
\end{quote}

\subsection*{You'll want to know: \mr{व्याज}}
\textbf{\mr{व्याज}} — \emph{vyāj} — \textquotedblleft{}interest on money\textquotedblright{}.

\begin{quote}
\mr{मी} \mr{दर} \mr{महिन्याला} \mr{कर्जावर} \mr{व्याज} \mr{भरतो}.
\end{quote}

\begin{quote}
\emph{mī dar mahinyālā karjāvar vyāj bharto.}
\end{quote}

\begin{quote}
Every month I pay interest on the loan.
\end{quote}

Neuter.

\mr{व्याज} \mr{मिळणे}, to earn interest; \mr{व्याज} \mr{भरणे}, to pay it.

\begin{cousinweb}
\textbf{\mr{व्याज}} (\emph{vyāj}) is the Marathi word for interest on money; its origin is not settled.
\end{cousinweb}

\begin{grammarlens}[title={What you've built}]
One more word for getting around.
\end{grammarlens}

\subsection*{Your turn}
\begin{itemize}
\raggedright
  \item \emph{Say it:} \emph{vyāj}
  \item \emph{Say it:} \emph{vyāj}, once more
  \item \emph{Say it:} \emph{mī dar mahinyālā karjāvar vyāj bharto.} — Every month I pay interest on the loan.
  \item \emph{Recall:} say the Marathi for a sum of money, then the Marathi for to deposit, then say \emph{vyāj} again
\end{itemize}

\begin{tcolorbox}[breakable,colback=teal!4,colframe=teal!35!black,title={Before you move on}]
What does \mr{व्याज} mean? (\textquotedblleft{}Interest on money\textquotedblright{}.) Say it once more.
\end{tcolorbox}
\section[\texorpdfstring{\mr{थकबाकी} (thakbākī)}{थकबाकी (thakbākī)}]{\mr{थकबाकी} (thakbākī) — arrears}

\label{lesson:MR-C427-thakbaki}



\begin{quote}
Before the new one: say the Marathi for to deposit, then the Marathi for interest on money.
\end{quote}

\subsection*{You'll want to know: \mr{थकबाकी}}
\textbf{\mr{थकबाकी}} — \emph{thakbākī} — \textquotedblleft{}arrears\textquotedblright{}.

\begin{quote}
\mr{त्याच्यावर} \mr{विजेच्या} \mr{बिलाची} \mr{थकबाकी} \mr{आहे}.
\end{quote}

\begin{quote}
\emph{tyācyāvar vijecyā bilācī thakbākī āhe.}
\end{quote}

\begin{quote}
He has arrears on his electricity bill.
\end{quote}

Feminine.

Money owed and overdue; the person owing takes -\mr{वर}.

\begin{cousinweb}
\textbf{\mr{थकबाकी}} (\emph{thakbākī}) is built from the stem of \textbf{\mr{थकणे}}, \textquotedblleft{}to tire, to fall behind\textquotedblright{}, and \textbf{\mr{बाकी}}, \textquotedblleft{}what remains\textquotedblright{}, a loan usually traced to Arabic \emph{bāqī}.
\end{cousinweb}

\begin{grammarlens}[title={What you've built}]
One more word for getting around.
\end{grammarlens}

\subsection*{Your turn}
\begin{itemize}
\raggedright
  \item \emph{Say it:} \emph{thakbākī}
  \item \emph{Say it:} \emph{thakbākī}, once more
  \item \emph{Say it:} \emph{tyācyāvar vijecyā bilācī thakbākī āhe.} — He has arrears on his electricity bill.
  \item \emph{Recall:} say the Marathi for to deposit, then the Marathi for interest on money, then say \emph{thakbākī} again
\end{itemize}

\begin{tcolorbox}[breakable,colback=teal!4,colframe=teal!35!black,title={Before you move on}]
What does \mr{थकबाकी} mean? (\textquotedblleft{}Arrears\textquotedblright{}.) Say it once more.
\end{tcolorbox}
\section[\texorpdfstring{\mr{उधार} (udhār)}{उधार (udhār)}]{\mr{उधार} (udhār) — on credit}

\label{lesson:MR-C427-udhar}



\begin{quote}
Before the new one: say the Marathi for interest on money, then the Marathi for arrears.
\end{quote}

\subsection*{You'll want to know: \mr{उधार}}
\textbf{\mr{उधार}} — \emph{udhār} — \textquotedblleft{}on credit\textquotedblright{}.

\begin{quote}
\mr{तो} \mr{वाण्याकडून} \mr{उधार} \mr{सामान} \mr{घेतो}.
\end{quote}

\begin{quote}
\emph{to vāṇyākaḍūn udhār sāmān gheto.}
\end{quote}

\begin{quote}
He buys groceries on credit from the grocer.
\end{quote}

\mr{उधार} \mr{घेणे}, to take on credit; \mr{उधार} \mr{देणे}, to give on credit.

\begin{cousinweb}
\textbf{\mr{उधार}} (\emph{udhār}) is inherited: CDIAL lists Marathi \emph{udhār}, \textquotedblleft{}buying and selling on credit\textquotedblright{}, under Sanskrit \emph{uddhāra}, \textquotedblleft{}raising; a loan\textquotedblright{} (CDIAL 2018).
\end{cousinweb}

\begin{grammarlens}[title={What you've built}]
One more word for getting around.
\end{grammarlens}

\subsection*{Your turn}
\begin{itemize}
\raggedright
  \item \emph{Say it:} \emph{udhār}
  \item \emph{Say it:} \emph{udhār}, once more
  \item \emph{Say it:} \emph{to vāṇyākaḍūn udhār sāmān gheto.} — He buys groceries on credit from the grocer.
  \item \emph{Recall:} say the Marathi for interest on money, then the Marathi for arrears, then say \emph{udhār} again
\end{itemize}

\begin{tcolorbox}[breakable,colback=teal!4,colframe=teal!35!black,title={Before you move on}]
What does \mr{उधार} mean? (\textquotedblleft{}On credit\textquotedblright{}.) Say it once more.
\end{tcolorbox}
\section[\texorpdfstring{\mr{बाकी} (bākī)}{बाकी (bākī)}]{\mr{बाकी} (bākī) — the remainder; the rest}

\label{lesson:MR-C427-baki}



\begin{quote}
Before the new one: say the Marathi for arrears, then the Marathi for on credit.
\end{quote}

\subsection*{You'll want to know: \mr{बाकी}}
\textbf{\mr{बाकी}} — \emph{bākī} — \textquotedblleft{}the remainder; the rest\textquotedblright{}.

\begin{quote}
\mr{मी} \mr{अर्धे} \mr{पैसे} \mr{दिले}, \mr{बाकी} \mr{उद्या} \mr{देईन}.
\end{quote}

\begin{quote}
\emph{mī ardhe paise dile, bākī udyā deīn.}
\end{quote}

\begin{quote}
I paid half the money; I will give the rest tomorrow.
\end{quote}

Feminine.

Also the balance of an account: \mr{बाकी} \mr{किती}?

\begin{cousinweb}
\textbf{\mr{बाकी}} (\emph{bākī}) is a loan through Persian, usually traced to Arabic \emph{bāqī}, \textquotedblleft{}remaining\textquotedblright{}.
\end{cousinweb}

\begin{grammarlens}[title={What you've built}]
One more word for getting around.
\end{grammarlens}

\subsection*{Your turn}
\begin{itemize}
\raggedright
  \item \emph{Say it:} \emph{bākī}
  \item \emph{Say it:} \emph{bākī}, once more
  \item \emph{Say it:} \emph{mī ardhe paise dile, bākī udyā deīn.} — I paid half the money; I will give the rest tomorrow.
  \item \emph{Recall:} say the Marathi for arrears, then the Marathi for on credit, then say \emph{bākī} again
\end{itemize}

\begin{tcolorbox}[breakable,colback=teal!4,colframe=teal!35!black,title={Before you move on}]
What does \mr{बाकी} mean? (\textquotedblleft{}The remainder; the rest\textquotedblright{}.) Say it once more.
\end{tcolorbox}
\section[\texorpdfstring{\mr{फेडणे} (pheḍṇe)}{फेडणे (pheḍṇe)}]{\mr{फेडणे} (pheḍṇe) — to repay}

\label{lesson:MR-C427-phedne}



\begin{quote}
Before the new one: say the Marathi for on credit, then the Marathi for the remainder.
\end{quote}

\subsection*{You'll want to know: \mr{फेडणे}}
\textbf{\mr{फेडणे}} — \emph{pheḍṇe} — \textquotedblleft{}to repay\textquotedblright{}.

\begin{quote}
\mr{त्याने} \mr{पाच} \mr{वर्षांत} \mr{कर्ज} \mr{फेडलं}.
\end{quote}

\begin{quote}
\emph{tyāne pāc varṣānt karja pheḍlaṁ.}
\end{quote}

\begin{quote}
He repaid the loan in five years.
\end{quote}

Past, with -\mr{ने} (he …): \mr{त्याने} \mr{फेडलं}, \emph{tyāne pheḍlaṁ}.

\mr{कर्ज} \mr{फेडणे}, to pay off a debt; the debt is the object.

\begin{cousinweb}
\textbf{\mr{फेडणे}} (\emph{pheḍṇe}) is inherited: CDIAL lists Marathi \emph{pheḍṇẽ}, \textquotedblleft{}to loosen\textquotedblright{}, under \emph{*sphiṭati}, \textquotedblleft{}gets loose\textquotedblright{} (CDIAL 13838); paying off a debt is its everyday sense today.
\end{cousinweb}

\begin{grammarlens}[title={What you've built}]
One more word for getting around.
\end{grammarlens}

\subsection*{Your turn}
\begin{itemize}
\raggedright
  \item \emph{Say it:} \emph{pheḍṇe}
  \item \emph{Say it:} \emph{pheḍṇe}, once more
  \item \emph{Say it:} \emph{tyāne pāc varṣānt karja pheḍlaṁ.} — He repaid the loan in five years.
  \item \emph{Recall:} say the Marathi for on credit, then the Marathi for the remainder, then say \emph{pheḍṇe} again
\end{itemize}

\begin{tcolorbox}[breakable,colback=teal!4,colframe=teal!35!black,title={Before you move on}]
What does \mr{फेडणे} mean? (\textquotedblleft{}To repay\textquotedblright{}.) Say it once more.
\end{tcolorbox}
\section[(dialogue)]{Looking back: from Embassies and Zoos to Interest and Arrears}

\label{lesson:MR-R427-review-420-427}



\begin{quote}
Say the Marathi for an embassy and a planetarium.
\end{quote}

\subsection*{Your turn}
Say these aloud:

\begin{itemize}
\raggedright
  \item an embassy, a planetarium, a zoo, a resident, the surroundings
  \item construction, a floor, a slum, a village gate, a boundary
  \item a vegetable market, noise, an uproar, bustling, bustle
  \item to go shopping, a deal, a price, to bargain, to buy
  \item a sale, a hawker, wholesale, a broker, a bill
  \item a pair of scales, to measure, genuine, fake, a copy
  \item cash, a cheque, small change, a sum of money, to deposit
  \item interest on money, arrears, on credit, the remainder, to repay
\end{itemize}

\begin{tcolorbox}[breakable,colback=teal!4,colframe=teal!35!black,title={Before you move on}]
An embassy? (\textbf{\mr{दूतावास}}, \emph{dūtāvās}.) A planetarium? (\textbf{\mr{तारांगण}}, \emph{tārāṅgaṇ}.) A zoo? (\textbf{\mr{प्राणिसंग्रहालय}}, \emph{prāṇisaṅgrahālay}.) A resident? (\textbf{\mr{रहिवासी}}, \emph{rahivāsī}.) The surroundings? (\textbf{\mr{परिसर}}, \emph{parisar}.) Construction? (\textbf{\mr{बांधकाम}}, \emph{bāndhkām}.) A floor? (\textbf{\mr{मजला}}, \emph{majlā}.) A slum? (\textbf{\mr{झोपडपट्टी}}, \emph{jhopaḍpaṭṭī}.) A village gate? (\textbf{\mr{वेस}}, \emph{ves}.) A boundary? (\textbf{\mr{हद्द}}, \emph{hadda}.) A vegetable market? (\textbf{\mr{मंडई}}, \emph{maṇḍaī}.) Noise? (\textbf{\mr{गोंगाट}}, \emph{goṅgāṭ}.) An uproar? (\textbf{\mr{कोलाहल}}, \emph{kolāhal}.) Bustling? (\textbf{\mr{गजबजलेला}}, \emph{gajbajlelā}.) Bustle? (\textbf{\mr{वर्दळ}}, \emph{vardaḷ}.) To go shopping? (\textbf{\mr{खरेदी} \mr{करणे}}, \emph{kharedī karṇe}.) A deal? (\textbf{\mr{सौदा}}, \emph{saudā}.) A price? (\textbf{\mr{किंमत}}, \emph{kimmat}.) To bargain? (\textbf{\mr{भाव} \mr{करणे}}, \emph{bhāv karṇe}.) To buy? (\textbf{\mr{विकत} \mr{घेणे}}, \emph{vikat gheṇe}.) A sale? (\textbf{\mr{विक्री}}, \emph{vikrī}.) A hawker? (\textbf{\mr{फेरीवाला}}, \emph{pherīvālā}.) Wholesale? (\textbf{\mr{घाऊक}}, \emph{ghāūk}.) A broker? (\textbf{\mr{दलाल}}, \emph{dalāl}.) A bill? (\textbf{\mr{बिल}}, \emph{bil}.) A pair of scales? (\textbf{\mr{तराजू}}, \emph{tarājū}.) To measure? (\textbf{\mr{मापणे}}, \emph{māpṇe}.) Genuine? (\textbf{\mr{अस्सल}}, \emph{assal}.) Fake? (\textbf{\mr{बनावट}}, \emph{banāvaṭ}.) A copy? (\textbf{\mr{नक्कल}}, \emph{nakkal}.) Cash? (\textbf{\mr{रोकड}}, \emph{rokaḍ}.) A cheque? (\textbf{\mr{धनादेश}}, \emph{dhanādeś}.) Small change? (\textbf{\mr{चिल्लर}}, \emph{cillar}.) A sum of money? (\textbf{\mr{रक्कम}}, \emph{rakkam}.) To deposit? (\textbf{\mr{जमा} \mr{करणे}}, \emph{jamā karṇe}.) Interest on money? (\textbf{\mr{व्याज}}, \emph{vyāj}.) Arrears? (\textbf{\mr{थकबाकी}}, \emph{thakbākī}.) On credit? (\textbf{\mr{उधार}}, \emph{udhār}.) The remainder? (\textbf{\mr{बाकी}}, \emph{bākī}.) To repay? (\textbf{\mr{फेडणे}}, \emph{pheḍṇe}.)
\end{tcolorbox}
